% Imported and adapted from the closed standalone d=4 manuscript.
% All labels are prefixed d4: to keep the combined paper map unambiguous.

The second main result evaluates the complete first-defect diagonal.

\begin{samepage}
\begin{theorem}[First defect diagonal]\label{d4:thm:main}
Let $T=T(z)\in\Q[[z]]$ be the unique series with constant coefficient $1$
such that $T=1+zT^3$, and put
\[
 D=(3-2T)(T^2-3T+3).
\]
Then
\begin{equation}\label{d4:eq:main-formula}
 \sum_{m\ge0}T_{103}(m+4,2m+4)z^m
 =\frac{T^3}{D}
   (9-6T+27T^2-35T^3+11T^4+2T^5-T^6).
\end{equation}
\end{theorem}
\end{samepage}

\section{Finite defects and the first defect diagonal}

We use an integer-coordinate realization of the hexagonal grid.  A cell is a
pair $(i,j)\in\mathbb Z^2$, with $\kappa=1-i-j$.  The $(a,b)$-benzel is
\begin{equation}\label{d4:eq:benzel}
 \Bz(a,b)=\left\{(i,j):
 1-a\le j-i,\ \kappa-j,\ i-\kappa\le b-1\right\}.
\end{equation}
This is equivalent to the geometric definition in~\cite{DLPY}.

Put $n=a-2$.  An \emph{owner anchor} is a pair $(q,r)$ satisfying
\[
 q-r\equiv n\pmod 3.
\]
It owns the three labelled cells
\[
 (q,r)_0,\qquad(q+1,r)_1,\qquad(q,r+1)_2.
\]
Every grid cell has exactly one owner and one label.  For an owner meeting the
benzel, define
\begin{equation}\label{d4:eq:uvw}
 u=\frac{n-q+r}{3},\qquad
 v=\frac{n-q-2r}{3},\qquad
 w=\frac{n+2q+r}{3}.
\end{equation}
Then $u,v,w\in\N$ and $u+v+w=n$.  Conversely,
\[
 q=w-u,\qquad r=u-v,
\]
so the owner is recovered from its simplex point.

The three cells owned by $(u,v,w)$ have difference triples
\begin{equation}\label{d4:eq:difference-table}
\begin{array}{c|ccc}
\text{label}&j-i&\kappa-j&i-\kappa\\ \midrule
0&3u-n&3v-n+1&3w-n-1\\
1&3u-n-1&3v-n&3w-n+1\\
2&3u-n+1&3v-n-1&3w-n.
\end{array}
\end{equation}
Substitution into~\eqref{d4:eq:benzel} gives the following truncation.

\begin{lemma}[Owner-domain truncation]\label{d4:lem:owner-domain}
Suppose $d=3k$ or $d=3k+1$, the benzel parameters are admissible, and
$2k\le n+1$.  The
owners meeting $\Bz(a,2a-d)$ are exactly
\[
 \Delta_n^{(k)}=\{(u,v,w)\in\N^3:u+v+w=n,\
     \max(u,v,w)\le n-k+1\}.
\]
Each of the three boundary segments $u=n-k+1$, $v=n-k+1$, and
$w=n-k+1$ contains $k$ deficient owners.  If $d=3k$, each deficient owner
contains two cells; if $d=3k+1$, it contains one cell.
\end{lemma}

\begin{proof}
For $d=3k$, label $0$ is present exactly when
\[
 u,w\le n-k+1,\qquad v\le n-k,
\]
and the other two conditions are obtained cyclically.  For $d=3k+1$, label
$0$ is present exactly when
\[
 u,v\le n-k,\qquad w\le n-k+1,
\]
again with cyclic analogues.  Taking the union over labels gives
$\Delta_n^{(k)}$.  On, for example, $u=n-k+1$, the remaining coordinates
are $(v,w)=(j,k-1-j)$ for $0\le j<k$, proving the boundary count and the
claimed multiplicities.
\end{proof}

Define the label potentials at owner $(q,r)$ by
\begin{equation}\label{d4:eq:potential}
 h_0(q,r)=q+r,\qquad h_1(q,r)=-q,\qquad h_2(q,r)=-r.
\end{equation}
Their sum at a full owner is zero.  A direct check on the four allowed
prototiles gives the local energy table
\begin{equation}\label{d4:eq:energy-table}
\begin{array}{c|cccc}
\text{tile class}&\substack{\text{in-phase}\\\text{stone}}
 &\substack{\text{wrong-phase}\\\text{stone}}
 &\substack{\text{two-owner}\\\text{bone}}
 &\substack{\text{three-owner}\\\text{bone}}\\ \midrule
\text{energy}&0&3&1&4.
\end{array}
\end{equation}
A right stone is \emph{in phase} when it is one owner triple; otherwise it is
wrong phase.  A bone meets either two or three owners.

\begin{theorem}[Finite-defect theorem]\label{d4:thm:defects}
Put $n=a-2$, and suppose $k\ge1$ and $2k\le n+1$.  Assume that
$\Bz(a,2a-d)$ has a type-$103$ tiling.  Let $s$ be the number of wrong-phase
stones and $f$ the number of three-owner bones.
\begin{enumerate}
\item If $d=3k$, then $s+f=\binom{k}{2}$.
\item If $d=3k+1$, then $s+f=\binom{k+1}{2}$.
\end{enumerate}
\end{theorem}

\begin{proof}
Along one deficient boundary, the sum of the potentials of its present cells
is
\[
 \sum_{j=0}^{k-1}(n-2k+2+j)
   =k\left(n-\frac{3(k-1)}2\right).
\]
The three cyclic boundary sums are equal.  Hence the energy of the entire
region is
\begin{equation}\label{d4:eq:region-energy}
 E_{a,d}=3k\left(a-\frac{3k+1}{2}\right).
\end{equation}
Using the benzel area and Conway--Lagarias invariant from~\cite{DLPY}, the
number $B$ of bones in any type-$103$ tiling is
\begin{equation}\label{d4:eq:bone-counts}
 B=\begin{cases}
 3k(a-2k),&d=3k,\\
 3k(a-2k-1),&d=3k+1.
 \end{cases}
\end{equation}
Here is the substitution explicitly.  In terms of $n=a-2$, the area formula
gives the following numbers $N$ of trihexes, while the invariant gives the
numbers $S$ of right stones:
\[
\begin{array}{c|cc}
 &N&S\\ \midrule
d=3k
 &\dfrac{n^2+3n+2-3k^2+k}{2}
 &\dfrac{n^2+3n+2+9k^2-6kn-11k}{2}\\[5pt]
d=3k+1
 &\dfrac{n^2+3n+2-3k^2-k}{2}
 &\dfrac{n^2+3n+2+9k^2-6kn-7k}{2}.
\end{array}
\]
Since every tile has three cells and no left stones are allowed, $B=N-S$;
this is exactly~\eqref{d4:eq:bone-counts}.  Finally,
summing~\eqref{d4:eq:energy-table} over the tiling gives
\[
 E_{a,d}=B+3(s+f).
\]
Substitution of~\eqref{d4:eq:region-energy} and~\eqref{d4:eq:bone-counts} yields
the two formulas.
\end{proof}

\begin{remark}
If $d=3k+2$, all three labels at an owner have the same presence condition
$\max(u,v,w)\le n-k$.  The resulting all-right-stone tiling is the unique
stones-and-bones tiling by~\cite[Theorem~1.1]{DLPY}.  This known residue class
is not part of the new defect count in Theorem~\ref{d4:thm:defects}.
\end{remark}

The offsets $d=0,1,2$ have the unique all-right-stone tilings described in
\cite[Theorem~1.1 and Remark~5.1]{DLPY}; in particular they have no defect.
For $d=3$, Theorem~\ref{d4:thm:defects} again gives zero defects.  Thus $d=4$
is the first diagonal with a nonzero defect count, and that count is exactly
one.  We now classify and enumerate this one-defect case.

\section{The one-defect path bijection}

Set
\[
 a=m+4,\qquad b=2m+4,\qquad n=m+2.
\]
The owner domain is the full simplex $\Delta_n$, except that each of its three
corner owners contains a single boundary cell.  Every tiling has one bad
tile, of one of five kinds: two wrong-stone phases or one of the three bone
orientations in its three-owner phase.

The five possibilities admit the following uniform parameterization.  Write
$p=(x,y,z)$ and let $c_\ell$ be the owner of the label-$\ell$ cell of the bad
tile.  The condition $p\in\Delta_r$ means $x+y+z=r$.
\begin{equation}\label{d4:eq:defect-cores}
\begin{array}{c|c|ccc}
\text{class}&\text{parameter}&c_0&c_1&c_2\\ \midrule
S_1&p\in\Delta_m
 &(x+1,y,z+1)&(x+1,y+1,z)&(x,y+1,z+1)\\
S_2&p\in\Delta_{m+1}
 &(x,y,z+1)&(x+1,y,z)&(x,y+1,z)\\
B_A&p\in\Delta_m
 &(x,y,z+2)&(x+1,y+1,z)&(x,y+1,z+1)\\
B_B&p\in\Delta_m
 &(x+1,y,z+1)&(x+1,y+1,z)&(x,y+2,z)\\
B_C&p\in\Delta_m
 &(x+1,y,z+1)&(x+2,y,z)&(x,y+1,z+1).
\end{array}
\end{equation}
Here $S_1,S_2$ are the two wrong-stone phases, and $B_A,B_B,B_C$
are the three three-owner bone classes.  These formulas follow by applying
\eqref{d4:eq:uvw} to the three cells of each translated prototile; in
particular, they are literal placement coordinates rather than a quotient or
orbit parameterization.

Delete the bad tile.  Every remaining bone is a two-owner bone; orient it
from the owner containing two of its cells toward the owner containing its
third cell.  Label the edge by that third cell's label.  We call the resulting
directed graph the \emph{owner graph}.

In owner-anchor coordinates, put
\[
 \alpha=(2,-1),\qquad \beta=(-1,-1),\qquad \gamma=(-1,2).
\]
A label-$0$ edge has step $\alpha$ or $\gamma$, a label-$1$ edge has step
$\beta$ or $\gamma$, and a label-$2$ edge has step $\alpha$ or $\beta$.
Equivalently, in simplex coordinates these steps are
\[
 \alpha=(-1,0,1),\qquad
 \beta=(0,1,-1),\qquad
 \gamma=(1,-1,0).
\]
An \emph{admissible label-$\ell$ path} is a directed path from $c_\ell$ to
the one-cell boundary owner of label $\ell$, using the corresponding two
steps and remaining in $\Delta_n$.  The path with no edges is allowed exactly
when its two endpoints coincide.

\begin{lemma}[Three paths]\label{d4:lem:three-paths}
The owner graph is the disjoint union of three directed paths.  For each
$\ell\in\{0,1,2\}$, the label-$\ell$ path begins at the owner of the
label-$\ell$ cell of the bad tile and ends at the one-cell boundary owner of
label $\ell$.  A path may be empty when its initial owner is already its
boundary endpoint.
\end{lemma}

\begin{proof}
At a full owner away from the defect, the local exact-cover alternatives are:
an in-phase stone and no incident edge; one incoming and one outgoing edge of
the same label; or three incoming edges, one of each label.  This is the
complete six-class check supplied by the two allowed bone phases in each
orientation.  Thus every owner incident to an edge is a pass-through vertex
unless it is a source or terminal.  At a defect owner distinct from its
labelled boundary owner, the bad tile covers the label-$\ell$ cell and one
outgoing label-$\ell$ edge covers the other two cells.  At a boundary owner
not already occupied by the bad tile, one incoming edge covers its unique
cell.  If a defect owner is already its matching boundary owner, these two
incidences cancel and give the allowed empty path.

The label potentials in~\eqref{d4:eq:potential} strictly increase along
directed edges, so directed cycles are impossible.  Remove the empty paths
from the following count.  For a connected component $C$, let $S_C$ be its
number of defect sources, $B_C$ its number of one-cell boundary terminals,
and $K_C$ its number of internal three-incoming terminals.  Every remaining
vertex has one incoming and one outgoing edge.  Summing outdegree minus
indegree over $C$ therefore gives
\[
 S_C=B_C+3K_C.
\]
Globally the numbers of nonempty defect sources and nonempty boundary
terminals are equal: both are three minus the number of empty paths.  Hence
$\sum_C K_C=0$, so no internal three-incoming terminal occurs.  With no
merger and no directed cycle, every nonempty component is a single directed
path with one defect source and one boundary terminal; an additional
component would have neither a source nor a terminal and would therefore
contain a directed cycle.  Local exact coverage at a pass-through owner
forces its incoming and outgoing edges to carry the same label.  The source
and terminal labels consequently agree, proving the stated three paths.
\end{proof}

The path steps have a useful monotonicity in simplex coordinates:
\begin{equation}\label{d4:eq:path-monotonicity}
\begin{array}{c|cc}
\text{label}&\text{nonincreasing}&\text{nondecreasing}\\ \midrule
0&v&w\\
1&w&u\\
2&u&v.
\end{array}
\end{equation}
For every one of the five defect classes, the three initial owners satisfy
strict cyclic separations
\begin{equation}\label{d4:eq:core-separation}
 w(c_1)<w(c_0),\qquad u(c_2)<u(c_1),\qquad v(c_0)<v(c_2).
\end{equation}
The inequalities are read directly from the five rows of
\eqref{d4:eq:defect-cores}.  For example, every row has
$w(c_1)<w(c_0)$.  A label-$0$ path can only increase $w$, whereas a label-$1$
path can only decrease it, so these two paths cannot meet.  The other two
inequalities give the same argument for labels $(1,2)$ using $u$ and labels
$(2,0)$ using $v$.  Hence independently chosen paths of successive labels
cannot share an owner.

Conversely, fix a bad tile and three admissible labelled paths.  Install the
bad tile, realize each path edge by its corresponding bone, and put the
in-phase right stone at every unused full owner.  At a nonterminal core owner,
the bad cell and first edge cover all present cells; for an empty path, the
bad tile covers the unique cell of the coincident core/boundary owner.  At a
pass-through owner, the incoming edge covers the path label and the outgoing
edge covers the other two; at a nonempty boundary endpoint, the incoming edge
covers the unique cell.  Thus the construction is a literal exact cover.  Its
local description also shows that it inverts deletion of the bad tile.

\begin{proposition}[Literal bijection]\label{d4:prop:bijection}
For every $m\ge0$, type-$103$ tilings of $\Bz(m+4,2m+4)$ are in bijection
with a choice of one of the five defect classes and three independently
admissible labelled paths from its core owners to the corresponding boundary
owners.
\end{proposition}

\section{Ballot enumeration}

For $p,q\in\N$, put
\begin{equation}\label{d4:eq:R}
 R(p,q)=\binom{p+2q}{q}-\binom{p+2q}{q-1},
\end{equation}
where the second term is zero when $q=0$.  By the reflection principle,
$R(p,q)$ counts walks starting at height $p$, with $q$ up-steps and $p+q$
down-steps, that never go below zero.  Equivalently,
\begin{equation}\label{d4:eq:R-coeff}
 R(p,q)=[x^q](1-x)(1+x)^{p+2q}.
\end{equation}

For a simplex point $(u,v,w)$, the three labelled path counts are cyclically
$R(u,v)$, $R(v,w)$, and $R(w,u)$.  Substituting the five core
parameterizations into these three factors gives the next theorem.

Indeed, read a path backwards from its boundary endpoint.  For label $0$, the
two reverse steps become the down- and up-steps of a ballot walk that starts
at height $u$, has $v$ up-steps and $u+v$ down-steps, and ends at height zero.
Thus its count is $R(u,v)$.  Cyclic permutation gives $R(v,w)$ and $R(w,u)$
for labels $1$ and $2$.  Applying this rule to~\eqref{d4:eq:defect-cores} gives
\[
\begin{array}{c|c}
\text{class}&\text{product of its three arm counts}\\ \midrule
S_1&R(x+1,y)R(y+1,z)R(z+1,x)\\
S_2&R(x,y)R(y,z)R(z,x)\\
B_A&R(x,y)R(y+1,z)R(z+1,x)\\
B_B&R(x+1,y)R(y+1,z)R(z,x)\\
B_C&R(x+1,y)R(y,z)R(z+1,x).
\end{array}
\]
The last two rows are cyclic reindexings of the $B_A$ sum.

\begin{theorem}[Exact cyclic ballot sum]\label{d4:thm:d4-sum}
For $m=a-4$, define
\begin{align}
 A_m&=\sum_{x+y+z=m+1}R(x,y)R(y,z)R(z,x),\label{d4:eq:A}\\
 C_m&=\sum_{x+y+z=m}R(x+1,y)R(y+1,z)R(z+1,x),\label{d4:eq:C}\\
 H_m&=\sum_{x+y+z=m}R(x,y)R(y+1,z)R(z+1,x).\label{d4:eq:H}
\end{align}
Then
\begin{equation}\label{d4:eq:ballot-sum}
 T_{103}(a,2a-4)=A_m+C_m+3H_m.
\end{equation}
\end{theorem}

\begin{proof}
The second wrong-stone phase is parameterized by
$x+y+z=m+1$ and gives~\eqref{d4:eq:A}; the first is parameterized by
$x+y+z=m$ and gives~\eqref{d4:eq:C}.  One three-owner bone orientation gives
\eqref{d4:eq:H}.  The other two orientations yield cyclic permutations of its
summand, hence the same total.  Proposition~\ref{d4:prop:bijection} makes the
three arm choices independent and proves that these five classes exhaust the
literal tilings without multiplicity.
\end{proof}

\section{Lagrange--Good evaluation}

We evaluate all coefficients at once.  Let $x_0,x_1,x_2$ be commuting
variables and read subscripts cyclically.  Define
\begin{equation}\label{d4:eq:phi}
 \phi_0=(1+x_0)^2(1+x_1),\quad
 \phi_1=(1+x_1)^2(1+x_2),\quad
 \phi_2=(1+x_2)^2(1+x_0),
\end{equation}
and $w_i=x_i/\phi_i$.  The specialized Good determinant is
\begin{equation}\label{d4:eq:determinant}
 \Delta(x)=
 \frac{1-x_0-x_1-x_2+x_0x_1+x_0x_2+x_1x_2-2x_0x_1x_2}
 {(1+x_0)(1+x_1)(1+x_2)}.
\end{equation}

The multivariate Lagrange--Good formula~\cite{Good,Hofbauer} applies to every
formal series
\[
 F\in\Q[[x_0,x_1,x_2]]
\]
and gives
\begin{equation}\label{d4:eq:good-expansion}
 F(x)=\sum_{\boldsymbol r\in\N^3}
 [x^{\boldsymbol r}]\,
 F(x)\phi(x)^{\boldsymbol r}\Delta(x)\,
 w(x)^{\boldsymbol r}.
\end{equation}
For completeness, the determinant in~\eqref{d4:eq:determinant} is obtained by
expanding
\[
 \det\left(\delta_{ij}-\frac{x_j}{\phi_i}
   \frac{\partial\phi_i}{\partial x_j}\right)_{0\le i,j\le2}.
\]

Introduce the three numerators
\begin{align*}
 N_A&=(1-x_0)(1-x_1)(1-x_2),\\
 N_C&=(1-x_0^2)(1-x_1^2)(1-x_2^2),\\
 N_H&=(1-x_0^2)(1-x_1)(1-x_2^2),
\end{align*}
and put
\[
 F_A=\frac{N_A}{\Delta},\qquad
 F_C=\frac{N_C}{\Delta},\qquad
 F_H=\frac{N_H}{\Delta}.
\]
For $F=\sum_{i,j,k\ge0}f_{i,j,k}x_0^ix_1^jx_2^k$, define its total-degree
diagonal by
\[
 \operatorname{Diag}F(z)=\sum_{r\ge0}
   \left(\sum_{i+j+k=r}f_{i,j,k}\right)z^r.
\]
Using~\eqref{d4:eq:R-coeff}, the coefficient in~\eqref{d4:eq:good-expansion} is,
respectively, the summand in~\eqref{d4:eq:A}, \eqref{d4:eq:C}, or~\eqref{d4:eq:H}.
Thus total-degree diagonalization of~\eqref{d4:eq:good-expansion} gives the three
component series.

To make this explicit, write
\[
 A_0(r)=\sum_{x+y+z=r}R(x,y)R(y,z)R(z,x)
\]
and let $A_0(z)=\sum_{r\ge0}A_0(r)z^r$, with analogous notation for
$C(z)$ and $H(z)$.  On the diagonal $x_0=x_1=x_2=s$, all three $w_i$ equal
\begin{equation}\label{d4:eq:q}
 q(s)=\frac{s}{(1+s)^3}.
\end{equation}
The determinant becomes
\begin{equation}\label{d4:eq:diag-delta}
 \Delta(s,s,s)=
 \frac{(1-2s)(1-s+s^2)}{(1+s)^3}.
\end{equation}
Consequently,
\begin{align}
 A_0(q(s))&=\frac{(1-s)^3(1+s)^3}
 {(1-2s)(1-s+s^2)},\label{d4:eq:A-good}\\
 C(q(s))&=\frac{(1-s)^3(1+s)^6}
 {(1-2s)(1-s+s^2)},\label{d4:eq:C-good}\\
 H(q(s))&=\frac{(1-s)^3(1+s)^5}
 {(1-2s)(1-s+s^2)}.\label{d4:eq:H-good}
\end{align}

Now let $T=1+zT^3$ and set $s=T-1$.  Equation~\eqref{d4:eq:q} gives
$q(s)=z$.  If
\[
 D=(3-2T)(T^2-3T+3),
\]
then~\eqref{d4:eq:A-good}--\eqref{d4:eq:H-good} become
\begin{equation}\label{d4:eq:three-components}
 A_0(z)=\frac{T^3(2-T)^3}{D},\qquad
 C(z)=\frac{T^6(2-T)^3}{D},\qquad
 H(z)=\frac{T^5(2-T)^3}{D}.
\end{equation}
Since $A_m=A_0(m+1)$ and $A_0(0)=1$,
\[
 \sum_{m\ge0}A_mz^m=\frac{A_0(z)-1}{z}.
\]
Combining this identity with~\eqref{d4:eq:ballot-sum} and
\eqref{d4:eq:three-components}, then using $z=(T-1)/T^3$, gives exactly the
rational function in Theorem~\ref{d4:thm:main}.
